\documentclass[11pt]{article}

\usepackage[utf8]{inputenc}
\usepackage[T1]{fontenc}
\usepackage{amsmath}
\usepackage{amssymb}
\usepackage{booktabs}
\usepackage{multirow}
\usepackage{siunitx}
\usepackage{xcolor}
\usepackage{geometry}
\usepackage{array}
\usepackage{caption}
\usepackage{hyperref}

\geometry{margin=1in}

\definecolor{colorPure}{RGB}{180,180,180}
\definecolor{colorRAG}{RGB}{100,149,200}
\definecolor{colorHybrid}{RGB}{31,73,125}

\sisetup{round-mode=places, round-precision=4}

\title{\textbf{Sensitivity Analysis for MAVT-Based MCDA}\\[0.4em]
       \large Weight Perturbation, Reranking, and Ranking Stability\\
       Across LLM Architecture Comparisons}
\author{Ahaan Nigam\\
        Downingtown East High School, Exton, PA}
\date{April 2026}

\begin{document}
\maketitle
\tableofcontents
\vspace{1em}
\hrule
\vspace{1em}

%% ─────────────────────────────────────────────────────────────────────────────
\section{Purpose and Scope}
\label{sec:purpose}

The primary finding of the main study is that the Hybrid architecture outperforms
RAG-Enhanced, which outperforms Pure Prompting, as measured by Kendall's $\tau$
against a physics-based ground truth.  That result is computed under a \emph{fixed}
set of MAVT criterion weights:
\[
  w_{\text{ene}} = 0.30, \quad
  w_{\text{env}} = 0.35, \quad
  w_{\text{com}} = 0.20, \quad
  w_{\text{pra}} = 0.15.
\]
A natural concern is whether the ordering
$\tau_{\text{Hybrid}} > \tau_{\text{RAG}} > \tau_{\text{Pure}}$
is an artifact of this particular weight vector or whether it holds
across a reasonable neighbourhood of weight choices.

This document describes the sensitivity analysis that tests exactly that question.
Section~\ref{sec:mavt} reviews the MAVT scoring model.
Section~\ref{sec:perturbation} defines the ten weight perturbation scenarios.
Section~\ref{sec:reranking} explains the reranking procedure.
Section~\ref{sec:example} works through a concrete HVAC scenario
step-by-step under every perturbation.
Section~\ref{sec:robustness} defines the robustness check reported in the output.

%% ─────────────────────────────────────────────────────────────────────────────
\section{MAVT Scoring Framework}
\label{sec:mavt}

\subsection{Score structure}

Each decision scenario presents $m$ alternatives $\{A_1, \ldots, A_m\}$
(typically $m = 3$) evaluated on four criteria
$\mathcal{C} = \{$energy cost, environmental, comfort, practicality$\}$.
A physics-based ground truth (GT) calculator assigns each
alternative a real-valued score $s_{ij} \in [0,10]$ on criterion $j$.
Each LLM architecture independently produces its own estimate $\hat{s}_{ij}$
of those same scores.

The Multi-Attribute Value Theory (MAVT) composite score for alternative $i$ is:
\begin{equation}
  V(A_i \mid \mathbf{w})
  \;=\;
  \sum_{j \in \mathcal{C}} w_j \, s_{ij},
  \qquad \sum_{j} w_j = 1, \quad w_j \geq 0.
  \label{eq:mavt}
\end{equation}
Alternatives are ranked within each scenario by $V$ in descending order
(rank 1 $=$ highest composite score $=$ best option).

\subsection{Key insight for sensitivity analysis}

Both the GT ranking \emph{and} every architecture's ranking are derived
from equation~\eqref{eq:mavt} applied to their respective score matrices.
When the weight vector $\mathbf{w}$ is perturbed, the ``correct'' ranking
shifts together with the architecture ranking, because both use the same
formula.  This means the comparison is always fair: we measure how well
each architecture tracks a moving target, not a fixed one.

Concretely, for a perturbed weight vector $\mathbf{w}'$:
\begin{itemize}
  \item GT rank of $A_i$ in scenario $k$
        $= \operatorname{rank}\!\left(V(A_i \mid \mathbf{w}')\right)$
        computed from GT criterion scores.
  \item Architecture rank of $A_i$ in scenario $k$
        $= \operatorname{rank}\!\left(\hat{V}(A_i \mid \mathbf{w}')\right)$
        computed from that architecture's criterion score estimates.
\end{itemize}
Kendall's $\tau$ is then computed by comparing these two rank sequences
using the same \texttt{compute\_ranking\_metrics()} function used in the
main evaluation, with no other changes.

%% ─────────────────────────────────────────────────────────────────────────────
\section{Weight Perturbation Methodology}
\label{sec:perturbation}

\subsection{Perturbation rule}

Starting from the baseline weight vector $\mathbf{w}^{(0)}$, we generate
perturbed vectors by changing a single criterion's weight by $\delta = \pm 0.05$
and redistributing the difference equally across the remaining three criteria:
\begin{equation}
  w_j' =
  \begin{cases}
    w_j^{(0)} + \delta & \text{if } j = \text{target criterion}, \\
    w_j^{(0)} - \delta/3 & \text{otherwise}.
  \end{cases}
  \label{eq:perturb}
\end{equation}
If any resulting weight falls below zero it is clipped to zero, and the
entire vector is renormalised to sum exactly to 1:
\begin{equation}
  w_j'' = \frac{\max(0,\, w_j')}{\sum_{k} \max(0,\, w_k')}.
  \label{eq:renorm}
\end{equation}
Because no baseline weight is close to zero, clipping never activates for
any of the eight single-criterion perturbations, so equation~\eqref{eq:perturb}
already satisfies the simplex constraint in this study.

\subsection{The ten scenarios}

Table~\ref{tab:weight-scenarios} lists all ten weight vectors used in the
sensitivity analysis.

\begin{table}[ht]
\centering
\caption{Ten weight perturbation scenarios. ``Ene''\,=\,energy cost,
  ``Env''\,=\,environmental, ``Com''\,=\,comfort, ``Pra''\,=\,practicality.
  The $\delta = +0.05$ and $\delta = -0.05$ rows are exact;
  the redistributed fractions $\tfrac{0.05}{3} \approx 0.0\overline{1}6$
  are rounded to four decimal places for display.}
\label{tab:weight-scenarios}
\begin{tabular}{lrrrr}
  \toprule
  Scenario     & $w_{\text{Ene}}$ & $w_{\text{Env}}$ & $w_{\text{Com}}$ & $w_{\text{Pra}}$ \\
  \midrule
  Baseline     & 0.3000 & 0.3500 & 0.2000 & 0.1500 \\
  \midrule
  ene $+0.05$  & \textbf{0.3500} & 0.3333 & 0.1833 & 0.1333 \\
  ene $-0.05$  & \textbf{0.2500} & 0.3667 & 0.2167 & 0.1667 \\
  \midrule
  env $+0.05$  & 0.2833 & \textbf{0.4000} & 0.1833 & 0.1333 \\
  env $-0.05$  & 0.3167 & \textbf{0.3000} & 0.2167 & 0.1667 \\
  \midrule
  com $+0.05$  & 0.2833 & 0.3333 & \textbf{0.2500} & 0.1333 \\
  com $-0.05$  & 0.3167 & 0.3667 & \textbf{0.1500} & 0.1667 \\
  \midrule
  pra $+0.05$  & 0.2833 & 0.3333 & 0.1833 & \textbf{0.2000} \\
  pra $-0.05$  & 0.3167 & 0.3667 & 0.2167 & \textbf{0.1000} \\
  \midrule
  Equal        & 0.2500 & 0.2500 & 0.2500 & 0.2500 \\
  \bottomrule
\end{tabular}
\end{table}

The perturbation of $\pm 0.05$ corresponds to roughly a 14--33\% relative
change per criterion and represents a meaningful, plausible disagreement
about the importance of any single criterion while keeping all weights
within a realistic range.  The equal-weight scenario is a classic stress
test in MCDA sensitivity analysis because it eliminates any prioritisation
assumptions entirely.

%% ─────────────────────────────────────────────────────────────────────────────
\section{Reranking Procedure}
\label{sec:reranking}

Algorithm~\ref{alg:rerank} describes the exact steps applied to the matched
data frame (output of \texttt{match\_scenarios()} after failure filtering)
for each of the ten weight scenarios.

\begin{figure}[ht]
\begin{center}
\fbox{
\begin{minipage}{0.88\textwidth}
\textbf{Algorithm 1: Rerank matched scenarios under weight vector $\mathbf{w}'$}
\medskip
\begin{enumerate}
  \item For each row $(k, i)$ (scenario $k$, alternative $i$) in the matched
        data frame, compute:
        \[
          V^{\text{GT}}_{ki} = \sum_{j \in \mathcal{C}} w'_j \cdot s^{\text{GT}}_{kij}
          \qquad
          \hat{V}^{\text{arch}}_{ki} = \sum_{j \in \mathcal{C}} w'_j \cdot \hat{s}^{\text{arch}}_{kij}
        \]
  \item Within each scenario $k$, assign GT ranks by sorting alternatives
        on $V^{\text{GT}}_{ki}$ in descending order; ties broken by
        \texttt{method='average'}:
        \[
          r^{\text{GT}}_{ki} = \operatorname{rank}_{\downarrow}(V^{\text{GT}}_{ki})
        \]
  \item Assign architecture ranks analogously from $\hat{V}^{\text{arch}}_{ki}$:
        \[
          r^{\text{arch}}_{ki} = \operatorname{rank}_{\downarrow}(\hat{V}^{\text{arch}}_{ki})
        \]
  \item Replace \texttt{gt\_rank} and \texttt{arch\_rank} columns with
        $r^{\text{GT}}_{ki}$ and $r^{\text{arch}}_{ki}$.  All criterion score
        columns are left untouched.
  \item Pass the updated data frame to \texttt{compute\_ranking\_metrics()},
        which computes per-scenario Kendall's $\tau$ and Spearman's $\rho$
        and averages across scenarios.
\end{enumerate}
\end{minipage}
}
\end{center}
\caption{Reranking procedure applied for each weight perturbation scenario.}
\label{alg:rerank}
\end{figure}

Two design choices are worth noting:
\begin{itemize}
  \item \textbf{Pre-computed ranks from CSV files are discarded.}
        The CSV output files store ranks computed under the baseline weights.
        Using those ranks for perturbed scenarios would give incorrect
        ``targets'' because the GT ranking changes with the weights.
        Algorithm~\ref{alg:rerank} regenerates ranks from raw criterion
        scores on every call.
  \item \textbf{Failure-filtered rows are excluded once, before the loop.}
        Any scenario containing the failure sentinel value (1928) is
        removed once at load time.  The set of evaluated scenarios is then
        identical across all ten weight perturbations, so differences in
        $\tau$ reflect only the weight change, not a change in sample.
\end{itemize}

%% ─────────────────────────────────────────────────────────────────────────────
\section{Worked Example: HVAC Scenario 0}
\label{sec:example}

\subsection{Scenario description}

We trace Scenario~0 from the HVAC ground truth through the full sensitivity
pipeline.

\medskip
\begin{center}
\begin{tabular}{ll}
  \toprule
  Field & Value \\
  \midrule
  Question     & \textit{``The house will be empty for 8 hours today.}\\
               & \textit{\phantom{``}What AC temperature should I set?''} \\
  Location     & Chester Springs, PA \\
  Outdoor temp & $88\,^{\circ}\text{F}$ \\
  Home profile & 1,800 sq ft condo, poor insulation, 3 occupants \\
  Alternatives & $78\,^{\circ}\text{F}$, $80\,^{\circ}\text{F}$, $83\,^{\circ}\text{F}$ \\
  \bottomrule
\end{tabular}
\end{center}
\medskip

\subsection{Ground truth criterion scores}

Table~\ref{tab:gt-scores} shows the raw criterion scores produced by the
physics-based ground truth calculator for each alternative.

\begin{table}[ht]
\centering
\caption{Ground truth criterion scores for HVAC Scenario~0.
  All scores are on a $[0,10]$ scale; higher is better.}
\label{tab:gt-scores}
\begin{tabular}{lrrrr}
  \toprule
  Alternative & $s_{\text{Ene}}$ & $s_{\text{Env}}$ & $s_{\text{Com}}$ & $s_{\text{Pra}}$ \\
  \midrule
  $78\,^{\circ}\text{F}$ ($A_1$) & 5.90 & 5.90 & 8.60 & 9.59 \\
  $80\,^{\circ}\text{F}$ ($A_2$) & 6.24 & 6.24 & 6.11 & 10.00 \\
  $83\,^{\circ}\text{F}$ ($A_3$) & 6.76 & 6.76 & 2.86 & 8.72  \\
  \bottomrule
\end{tabular}
\end{table}

Note the trade-off structure: setting the AC higher ($83\,^{\circ}\text{F}$)
saves more energy and reduces emissions (higher energy cost and environmental
scores) but sharply penalises comfort ($s_{\text{Com}} = 2.86$).
The $78\,^{\circ}\text{F}$ setting sacrifices some energy efficiency to
maintain comfort.

\subsection{Baseline MAVT computation}

Under the baseline weights $(w_{\text{Ene}}, w_{\text{Env}},
w_{\text{Com}}, w_{\text{Pra}}) = (0.30, 0.35, 0.20, 0.15)$:

\begin{align}
  V(A_1) &= 0.30(5.90) + 0.35(5.90) + 0.20(8.60) + 0.15(9.59) \nonumber\\
         &= 1.770 + 2.065 + 1.720 + 1.439 = \mathbf{6.994} \quad (\text{rank } 1) \label{eq:a1} \\[4pt]
  V(A_2) &= 0.30(6.24) + 0.35(6.24) + 0.20(6.11) + 0.15(10.00) \nonumber\\
         &= 1.872 + 2.184 + 1.222 + 1.500 = \mathbf{6.778} \quad (\text{rank } 2) \label{eq:a2} \\[4pt]
  V(A_3) &= 0.30(6.76) + 0.35(6.76) + 0.20(2.86) + 0.15(8.72) \nonumber\\
         &= 2.028 + 2.366 + 0.572 + 1.308 = \mathbf{6.274} \quad (\text{rank } 3) \label{eq:a3}
\end{align}

The baseline GT ranking is $A_1 \succ A_2 \succ A_3$, i.e.\
$78\,^{\circ}\text{F}$ is the recommended setting.

\subsection{Sensitivity of the GT ranking to weight perturbations}

Table~\ref{tab:sensitivity-gt} shows the MAVT scores and resulting ranking
under all ten weight scenarios.

\begin{table}[ht]
\centering
\caption{GT MAVT scores and ranking for HVAC Scenario~0 under each weight
  perturbation.  Scores rounded to four decimal places.
  The margin column is $V(A_1) - V(A_2)$.}
\label{tab:sensitivity-gt}
\begin{tabular}{lrrrrll}
  \toprule
  Scenario    & $V(A_1)$ & $V(A_2)$ & $V(A_3)$ & Margin$(A_1{-}A_2)$ & Ranking \\
  \midrule
  Baseline    & 6.9935 & 6.7780 & 6.2740 & +0.2155 & $A_1 \succ A_2 \succ A_3$ \\
  ene $+0.05$ & 6.8870 & 6.7175 & 6.3063 & +0.1695 & $A_1 \succ A_2 \succ A_3$ \\
  ene $-0.05$ & 7.1000 & 6.8385 & 6.2417 & +0.2615 & $A_1 \succ A_2 \succ A_3$ \\
  env $+0.05$ & 6.8870 & 6.7175 & 6.3063 & +0.1695 & $A_1 \succ A_2 \succ A_3$ \\
  env $-0.05$ & 7.1000 & 6.8385 & 6.2417 & +0.2615 & $A_1 \succ A_2 \succ A_3$ \\
  com $+0.05$ & 7.0670 & 6.7088 & 6.0463 & +0.3582 & $A_1 \succ A_2 \succ A_3$ \\
  com $-0.05$ & 6.9200 & 6.8472 & 6.5017 & +0.0728 & $A_1 \succ A_2 \succ A_3$ \\
  pra $+0.05$ & 7.1330 & 6.9682 & 6.4370 & +0.1648 & $A_1 \succ A_2 \succ A_3$ \\
  pra $-0.05$ & 6.8540 & 6.5878 & 6.1110 & +0.2662 & $A_1 \succ A_2 \succ A_3$ \\
  Equal       & 7.4975 & 7.1475 & 6.2750 & +0.3500 & $A_1 \succ A_2 \succ A_3$ \\
  \bottomrule
\end{tabular}
\end{table}

\paragraph{Observation.}
The GT ranking $A_1 \succ A_2 \succ A_3$ is preserved under all ten weight
scenarios.  The margin between $A_1$ and $A_2$ narrows to a minimum of $0.073$
under com~$-0.05$ (where comfort weight drops to 0.15, weakening $A_1$'s
advantage in that criterion) but never changes sign.  Increasing the comfort
weight (com~$+0.05$) widens the margin to $0.358$, as $A_1$'s large comfort
score ($8.60$) gains influence.  The equal-weight scenario produces the
widest $A_1$--$A_2$ gap ($0.350$) because practicality, where $A_2$ scores
$10.00$ versus $A_1$'s $9.59$, is diluted to just $0.25$.

\subsection{Why the trade-off structure drives stability}

The key reason this scenario is stable is the trade-off structure visible in
Table~\ref{tab:gt-scores}: $A_1$ and $A_2$ score identically on energy cost
and environmental criteria (both quantitative; $A_1$: 5.90, $A_2$: 6.24).
$A_1$'s lead therefore comes almost entirely from comfort ($+2.49$ over $A_2$)
partially offset by practicality ($-0.41$).  For the ranking to flip,
$w_{\text{Com}}$ would need to drop far enough to erase a $2.49$-point comfort
advantage.  At the minimum tested value of $w_{\text{Com}} = 0.15$ (com~$-0.05$
scenario), that advantage contributes $0.15 \times 2.49 = 0.374$ score points
to $V(A_1)$, still more than enough to maintain first place.

This structural analysis illustrates the logic behind the full sensitivity
study: scenarios where criteria are closely scored and a single criterion is
the deciding factor are most likely to produce rank reversals under
perturbation.  Scenarios where the winner leads broadly across multiple
criteria are inherently more robust.

%% ─────────────────────────────────────────────────────────────────────────────
\section{Robustness Check}
\label{sec:robustness}

\subsection{Definition}

The sensitivity analysis reports the following robustness check for each of
the ten weight scenarios:
\begin{center}
  \textit{Is $\tau_{\text{Hybrid}} > \tau_{\text{RAG}} > \tau_{\text{Pure}}$
  in this scenario?}
\end{center}
A scenario is counted as ``preserved'' if the strict ordering holds.
The final summary line reports:
\[
  \frac{\text{number of scenarios where order is preserved}}{10}
\]

\subsection{Why strict inequality is used}

Average Kendall's $\tau$ is computed across all matched and non-failed scenarios
(typically 80--100 for this dataset).  Two architectures producing the same
averaged $\tau$ to four decimal places would require an extraordinary tie;
we therefore use strict inequality and expect it to hold in practice.

\subsection{Possible outcomes and their interpretation}

\begin{description}
  \item[$10/10$ preserved.]
    The architecture ordering is fully robust to reasonable weight
    perturbations.  The relative ranking of architectures does not depend
    on the specific weight choice made for the main study.
  \item[$7$--$9/10$ preserved.]
    The ordering holds under most perturbations.  The failing scenarios
    should be inspected to determine which weight direction breaks the
    ordering and why.
  \item[$\leq 6/10$ preserved.]
    The architecture ordering is sensitive to weight choice.  Results
    should be reported conditionally and the weight-specific performance
    differences highlighted.
\end{description}

\subsection{Scope of the sensitivity claim}

The analysis makes one deliberate simplification: all criterion weights are
treated as independent parameters for perturbation, and only single-criterion
perturbations of magnitude $\delta = 0.05$ are tested.  It does not cover
simultaneous multi-criterion perturbations or larger perturbations.
For the main study findings, the $\pm 0.05$ range is considered the
policy-relevant neighbourhood: it captures reasonable expert disagreement
about relative criterion importance without changing the fundamental
problem structure.

%% ─────────────────────────────────────────────────────────────────────────────
\section*{Summary}

\begin{enumerate}
  \item The sensitivity analysis perturbs the MAVT weight vector by $\pm 0.05$
        on each of four criteria in turn, plus an equal-weight baseline,
        yielding ten total scenarios.
  \item For each weight scenario, GT and architecture rankings are independently
        recomputed from raw criterion scores using equation~\eqref{eq:mavt}.
        Pre-computed ranks from CSV files are not used.
  \item Kendall's $\tau$ is computed between the reranked GT and each
        architecture, averaged across all successfully matched scenarios.
  \item HVAC Scenario~0 illustrates that a structurally dominant alternative
        (one that leads across multiple criteria) produces a ranking stable
        to all tested perturbations, with the winning margin narrowing or
        widening but never reversing.
  \item The robustness check counts how many of the ten scenarios preserve
        the ordering $\tau_{\text{Hybrid}} > \tau_{\text{RAG}} > \tau_{\text{Pure}}$.
        A full $10/10$ count confirms that the main result does not depend on
        the specific baseline weight vector.
\end{enumerate}

\end{document}
